\section{More Experimental Results}
\label{appendix:more_exp}


\subsection{More Comparison Experiments}
\label{appendix:more_exp_comparison}




We further compare MemSkill with additional baselines on ALFWorld, including LightMem~\citep{fang2025lightmem}, AWM~\citep{wang2024agent}, and Expel~\citep{zhao2024expel}. As shown in Table~\ref{tab:more_exp_comparison}, MemSkill consistently achieves the best performance across both base models and evaluation splits. With LLaMA, MemSkill obtains an average success rate of 80.36, outperforming the strongest baseline LightMem by 5.53 points. The gain is especially clear on ALF-Unseen, where MemSkill improves the success rate from 75.37 to 83.58 while reducing the average number of steps from 20.69 to 16.63.

The advantage remains under transfer evaluation with Qwen, where MemSkill is not trained using this base model or dataset. MemSkill achieves an average success rate of 81.29, surpassing the strongest baseline Expel by 7.61 points. It also requires fewer interaction steps on both ALF-Seen and ALF-Unseen. These results suggest that the learned memory skills can capture reusable experience patterns for embodied decision-making, leading to both higher task success and more efficient execution across seen and unseen environments.



\begin{table*}[h]
    \footnotesize
    \centering
    \setlength{\tabcolsep}{5.5pt}
    \caption{\textbf{More comparison results on ALFWorld.}}
    \label{tab:more_exp_comparison}

    \begin{tabular}{>{\centering\arraybackslash}m{0.5cm} l cc cc c}
        \toprule
        \multirow{3}{*}{\textbf{Model}} &
        \multirow{3}{*}{\textbf{Methods}} &
        \multicolumn{5}{c}{\textbf{Embodied Interactive Tasks}} \\
        \cmidrule(lr){3-7}
        & &
        \multicolumn{2}{c}{\textbf{ALF-Seen$^\dagger$}} &
        \multicolumn{2}{c}{\textbf{ALF-Unseen$^\dagger$}} &
        \multicolumn{1}{c}{\textbf{Avg.}} \\
        \cmidrule(lr){3-4} \cmidrule(lr){5-6} \cmidrule(lr){7-7}
        & &
        \textbf{SR} & \textbf{\#Stps$\downarrow$} &
        \textbf{SR} & \textbf{\#Stps$\downarrow$} &
        \textbf{SR} \\
        \midrule

        % ===================== Base model block 1 =====================
        \multirow{4}{*}{\rotatebox[origin=c]{90}{\shortstack{\textbf{LLaMA}}}} &
        LightMem & 74.29 & 21.69 & 75.37 & 20.69 & 74.83 \\
        & AWM   & 66.43 & 23.25 & 68.66 & 22.25 & 67.55 \\
        & Expel   & 67.14 & 23.26 & 66.42 & 22.40 & 66.78 \\
        & \cellcolor{hl-llama}\textbf{MemSkill} &
        \cellcolor{hl-llama}\textbf{77.14} & \cellcolor{hl-llama}\textbf{18.91} &
        \cellcolor{hl-llama}\textbf{83.58} & \cellcolor{hl-llama}\textbf{16.63} &
        \cellcolor{hl-llama}\textbf{80.36} \\
        \midrule

        % ===================== Base model block 2 =====================
        \multirow{4}{*}{\rotatebox[origin=c]{90}{\shortstack{$^\blacktriangle$\textbf{Qwen}}}} &
        LightMem & 70.71 & 20.48 & 58.21 & 25.62 & 64.46 \\
        & AWM   & 74.29 & 19.03 & 61.19 & 24.67 & 67.74 \\
        & Expel   & 75.71 & 18.81 & 71.64 & 20.98 & 73.68 \\
        & \cellcolor{hl-qwen}\textbf{MemSkill} &
        \cellcolor{hl-qwen}\textbf{85.71} & \cellcolor{hl-qwen}\textbf{13.84} &
        \cellcolor{hl-qwen}\textbf{76.87} & \cellcolor{hl-qwen}\textbf{18.16} &
        \cellcolor{hl-qwen}\textbf{81.29} \\
        \bottomrule

        \multicolumn{7}{l}{\small \textbf{Bold} indicates the best score within each base model block.} \\
        \multicolumn{7}{l}{\small $^\blacktriangle$ indicates no training using this base model or dataset.} \\
        \multicolumn{7}{l}{\small $^\dagger$ indicates evaluation with in-context demonstrations.} \\
    \end{tabular}
\end{table*}







\begin{table*}[h]
    \footnotesize
    \centering
    \setlength{\tabcolsep}{5.5pt}
    \caption{\textbf{More comparison results on ALFWorld \textit{without} in-context demonstrations.}}
    \label{tab:more_exp_alfworld}

    \begin{tabular}{>{\centering\arraybackslash}m{0.5cm} l cc cc c}
        \toprule
        \multirow{3}{*}{\textbf{Model}} &
        \multirow{3}{*}{\textbf{Methods}} &
        \multicolumn{5}{c}{\textbf{Embodied Interactive Tasks}} \\
        \cmidrule(lr){3-7}
        & &
        \multicolumn{2}{c}{\textbf{ALF-Seen}} &
        \multicolumn{2}{c}{\textbf{ALF-Unseen}} &
        \multicolumn{1}{c}{\textbf{Avg.}} \\
        \cmidrule(lr){3-4} \cmidrule(lr){5-6} \cmidrule(lr){7-7}
        & &
        \textbf{SR} & \textbf{\#Stps$\downarrow$} &
        \textbf{SR} & \textbf{\#Stps$\downarrow$} &
        \textbf{SR} \\
        \midrule

        % ===================== Base model block 1 =====================
        \multirow{12}{*}{\rotatebox[origin=c]{90}{\shortstack{\textbf{LLaMA3.3}\\\textbf{70B-Instruct}}}} &
        No-Memory   & 17.14 & 43.74 & 20.15 & 42.99 & 18.65 \\
        & CoN   & 40.71 & 33.44 & 30.60 & 37.66 & 35.66 \\
        & ReadAgent   & 32.86 & 37.09 & 38.06 & 34.78 & 35.46 \\
        & MemoryBank   & 25.00 & 39.96 & 32.84 & 36.54 & 28.92 \\
        & A-MEM   & 24.29 & 40.51 & 28.36 & 38.83 & 26.33 \\
        & Mem0   & 32.86 & 36.47 & 32.09 & 37.32 & 32.48 \\
        & LangMem   & 37.86 & 34.39 & 35.07 & 35.70 & 36.47 \\
        & MemoryOS   & 15.71 & 43.74 & 14.18 & 44.54 & 14.95 \\
        & LightMem   & 45.71 & 31.76 & 46.27 & 30.66 & 45.99 \\
        & AWM   & 35.71 & 35.87 & 39.55 & 34.93 & 37.63 \\
        & Expel   & 42.14 & 34.61 & 43.28 & 33.25 & 42.71 \\
        & \cellcolor{hl-llama}\textbf{MemSkill} &
        \cellcolor{hl-llama}\textbf{55.71} & \cellcolor{hl-llama}\textbf{26.87} &
        \cellcolor{hl-llama}\textbf{59.70} & \cellcolor{hl-llama}\textbf{25.88} &
        \cellcolor{hl-llama}\textbf{57.71} \\
        \midrule

        % ===================== Base model block 2 =====================
        \multirow{12}{*}{\rotatebox[origin=c]{90}{\shortstack{$^\blacktriangle$\textbf{Qwen3-Next}\\\textbf{80B-A3B-Instruct}}}} &
        No-Memory   & 18.57 & 42.48 & 26.12 & 39.35 & 22.35 \\
        & CoN   & 57.86 & 25.81 & 53.73 & 28.40 & 55.80 \\
        & ReadAgent   & 53.57 & 27.88 & 54.48 & 27.41 & 54.03 \\
        & MemoryBank   & 37.86 & 35.15 & 38.06 & 34.99 & 37.96 \\
        & A-MEM   & 25.00 & 40.28 & 29.10 & 39.04 & 27.05 \\
        & Mem0   & 38.57 & 33.64 & 41.04 & 33.16 & 39.81 \\
        & LangMem   & 37.14 & 34.42 & 31.34 & 37.17 & 34.24 \\
        & MemoryOS   & 19.29 & 42.43 & 18.66 & 42.95 & 18.98 \\
        & LightMem   & 30.71 & 36.81 & 26.87 & 38.96 & 28.79 \\
        & AWM   & 47.86 & 31.24 & 47.01 & 31.69 & 47.44 \\
        & Expel   & 54.29 & 28.90 & 55.22 & 29.57 & 54.76 \\
        & \cellcolor{hl-qwen}\textbf{MemSkill} &
        \cellcolor{hl-qwen}\textbf{65.71} & \cellcolor{hl-qwen}\textbf{22.49} &
        \cellcolor{hl-qwen}\textbf{61.94} & \cellcolor{hl-qwen}\textbf{24.59} &
        \cellcolor{hl-qwen}\textbf{63.83} \\
        \bottomrule

        \multicolumn{7}{l}{\small \textbf{Bold} indicates the best score within each base model block.} \\
        \multicolumn{7}{l}{\small $^\blacktriangle$ indicates no training using this base model or dataset.} \\
    \end{tabular}
\end{table*}


Table~\ref{tab:more_exp_alfworld} further reports ALFWorld results \textit{without} in-context demonstrations. This setting is complementary to the main-table evaluation with demonstrations, but is more controlled for studying memory itself. Since in-context demonstrations can also act as an external form of memory, including them may confound the gains brought by each method's constructed memory. Therefore, this setting decouples demonstrations from memory construction and isolates the contribution of the learned memory mechanism.

Under this stricter setting, MemSkill still consistently outperforms all baselines across both base models and both evaluation splits. With LLaMA, MemSkill achieves an average success rate of 57.71, outperforming the strongest baseline LightMem by 11.72 points. The improvement is larger on ALF-Unseen, where MemSkill improves the success rate from 46.27 to 59.70, suggesting stronger generalization to unseen environments. MemSkill also requires fewer interaction steps than all baselines, reducing the average steps on ALF-Seen and ALF-Unseen to 26.87 and 25.88, respectively.

The same trend holds in the transfer setting with Qwen, where MemSkill is not trained using this base model or dataset. MemSkill achieves an average success rate of 63.83, improving over the strongest baseline by 8.03 points. It also obtains the lowest number of steps on both ALF-Seen and ALF-Unseen. These results show that the gains of MemSkill do not rely on in-context demonstrations, and that the learned memory skills provide reusable task experience that improves both success rate and execution efficiency.











\subsection{More Results on Appworld}
\label{appendix:more_exp_appworld}



\begin{table*}[t]
    \footnotesize
    \centering
    \setlength{\tabcolsep}{5pt}
    \caption{\textbf{More comparison results on AppWorld \textit{with} and \textit{without} in-context demonstrations.}}
    \label{tab:more_exp_appworld}

    \begin{tabular}{>{\centering\arraybackslash}m{0.5cm} l cc cc c | cc cc c}
        \toprule
        \multirow{3}{*}{\textbf{Model}} &
        \multirow{3}{*}{\textbf{Methods}} &
        \multicolumn{10}{c}{\textbf{Interactive Tool-Use Benchmark}} \\
        \cmidrule(lr){3-7} \cmidrule(lr){8-12}
        & &
        \multicolumn{2}{c}{\textbf{Test-N}} &
        \multicolumn{2}{c}{\textbf{Test-C}} &
        \multicolumn{1}{c|}{\textbf{Avg.}} &
        \multicolumn{2}{c}{\textbf{Test-N$^\dagger$}} &
        \multicolumn{2}{c}{\textbf{Test-C$^\dagger$}} &
        \multicolumn{1}{c}{\textbf{Avg.}} \\
        \cmidrule(lr){3-4} \cmidrule(lr){5-6} \cmidrule(lr){7-7}
        \cmidrule(lr){8-9} \cmidrule(lr){10-11} \cmidrule(lr){12-12}
        & &
        \textbf{PR} & \textbf{\#Steps$\downarrow$} &
        \textbf{PR} & \textbf{\#Steps$\downarrow$} &
        \textbf{PR} &
        \textbf{PR} & \textbf{\#Steps$\downarrow$} &
        \textbf{PR} & \textbf{\#Steps$\downarrow$} &
        \textbf{PR} \\
        \midrule

        % ===================== Base model block 1 =====================
        \multirow{10}{*}{\rotatebox[origin=c]{90}{\shortstack{\textbf{LLaMA3.3}\\\textbf{70B-Instruct}}}} &
        No-Memory   & 22.80 & 36.71 & 19.10 & 37.92 & 20.95 & 54.03 & 14.21 & 38.78 & 19.12 & 46.41 \\
        & ReadAgent  & 23.43 & 34.82 & 20.49 & 31.36 & 21.96 & 53.76 & 14.54 & 38.67 & 19.24 & 46.22 \\
        & MemoryBank & 23.28 & 35.66 & 19.96 & 33.34 & 21.62 & 53.66 & 14.89 & 38.46 & 18.84 & 46.06 \\
        & A-MEM      & 22.13 & 38.46 & 17.60 & 38.46 & 19.87 & 52.85 & 14.31 & 38.16 & 18.78 & 45.51 \\
        & Mem0       & 27.02 & 27.04 & 21.62 & 30.64 & 24.32 & 54.32 & 14.14 & 39.73 & 18.55 & 47.03 \\
        & MemoryOS   & 23.24 & 38.26 & 19.17 & 36.17 & 21.21 & 53.14 & 13.62 & 38.52 & 17.83 & 45.83 \\
        & LightMem   & 25.42 & 30.73 & 21.32 & 30.33 & 23.37 & 51.25 & 14.10 & 39.12 & 18.50 & 45.19 \\
        & AWM        & 29.00 & \textbf{17.73} & 21.83 & \textbf{20.48} & 25.42 & 54.68 & 14.26 & 39.54 & 18.64 & 47.11 \\
        & Expel      & 25.29 & 29.96 & 21.20 & 30.63 & 23.25 & 52.88 & 13.56 & 39.49 & 17.91 & 46.19 \\
        & \cellcolor{hl-llama}\textbf{MemSkill} &
        \cellcolor{hl-llama}\textbf{31.12} & \cellcolor{hl-llama}29.38 &
        \cellcolor{hl-llama}\textbf{22.29} & \cellcolor{hl-llama}32.17 &
        \cellcolor{hl-llama}\textbf{26.71} &
        \cellcolor{hl-llama}\textbf{54.84} & \cellcolor{hl-llama}\textbf{13.25} &
        \cellcolor{hl-llama}\textbf{39.79} & \cellcolor{hl-llama}\textbf{17.59} &
        \cellcolor{hl-llama}\textbf{47.32} \\
        \midrule

        % ===================== Base model block 2 =====================
        \multirow{10}{*}{\rotatebox[origin=c]{90}{\shortstack{$^\blacktriangle$\textbf{Qwen3-Next}\\\textbf{80B-A3B-Instruct}}}} &
        No-Memory   & 22.41 & 38.39 & 19.99 & 36.78 & 21.20 & 47.68 & 16.81 & 38.42 & 19.70 & 43.05 \\
        & ReadAgent  & 26.29 & 34.77 & 20.36 & 35.07 & 23.33 & 46.83 & 17.59 & 37.65 & 21.12 & 42.24 \\
        & MemoryBank & 24.41 & 36.94 & 20.28 & 36.19 & 22.35 & 46.50 & 16.98 & 37.18 & 22.03 & 41.84 \\
        & A-MEM      & 23.03 & 37.17 & 18.39 & 37.43 & 20.71 & 45.44 & 17.49 & 38.14 & 21.40 & 41.79 \\
        & Mem0       & 40.42 & \textbf{18.48} & 22.16 & 22.97 & 31.29 & 50.19 & 15.11 & 39.64 & 20.33 & 44.92 \\
        & MemoryOS   & 25.00 & 35.98 & 21.46 & 35.20 & 23.23 & 50.98 & 16.11 & 38.03 & 21.32 & 44.51 \\
        & LightMem   & 36.90 & 23.71 & 23.44 & 27.05 & 30.17 & 48.91 & 17.23 & 40.85 & 21.14 & 44.88 \\
        & AWM        & 27.02 & 23.90 & 22.05 & 22.63 & 24.54 & 53.23 & 16.77 & 41.26 & 22.95 & 47.25 \\
        & Expel      & 25.00 & 22.20 & 20.40 & \textbf{21.00} & 22.70 & 49.88 & 16.86 & 39.17 & 20.06 & 44.53 \\
        & \cellcolor{hl-qwen}\textbf{MemSkill} &
        \cellcolor{hl-qwen}\textbf{43.05} & \cellcolor{hl-qwen}26.11 &
        \cellcolor{hl-qwen}\textbf{25.41} & \cellcolor{hl-qwen}31.39 &
        \cellcolor{hl-qwen}\textbf{34.23} &
        \cellcolor{hl-qwen}\textbf{54.53} & \cellcolor{hl-qwen}\textbf{15.05} &
        \cellcolor{hl-qwen}\textbf{42.30} & \cellcolor{hl-qwen}\textbf{19.23} &
        \cellcolor{hl-qwen}\textbf{48.42} \\
        \bottomrule

        \multicolumn{12}{l}{\small \textbf{Bold} indicates the best score within each base model block.} \\
        \multicolumn{12}{l}{\small $^\blacktriangle$ indicates no training using this base model or dataset.} \\
        \multicolumn{12}{l}{\small $^\dagger$ indicates evaluation with in-context demonstrations.} \\
    \end{tabular}
\end{table*}



To further evaluate whether memory skills generalize beyond conversational and embodied household settings, we extend our study to AppWorld~\citep{trivedi2024appworld}, a more challenging interactive tool-use benchmark. We report results on both Test-Normal (Test-N) and Test-Challenge (Test-C), using Pass Rate (PR) and the average number of execution steps as evaluation metrics. Our default setting follows the stronger evaluation protocol with in-context demonstrations. In addition, we also report a controlled variant without demonstrations, which helps isolate the contribution of each method's constructed memory from the task guidance provided by demonstrations.

As shown in Table~\ref{tab:more_exp_appworld}, MemSkill achieves the best average PR across both base models under the default demonstration-based setting. With LLaMA, the No-Memory baseline already performs strongly, indicating that in-context demonstrations provide substantial task-specific guidance and make this setting partially saturated. As a result, the performance gaps among different memory methods become relatively small. Nevertheless, MemSkill still achieves the best overall pass rate and uses the fewest execution steps across both splits, suggesting that the learned memory skills can still improve efficiency even when demonstrations already provide strong external guidance.

The benefit of MemSkill is more evident with Qwen, where the demonstration-based setting leaves more room for memory-based improvement. In this case, MemSkill consistently outperforms prior memory methods and also improves over the No-Memory baseline. It further achieves the lowest execution steps on both Test-N and Test-C, indicating that MemSkill can effectively complement in-context demonstrations and support more efficient tool-use behavior.

The controlled setting without in-context demonstrations further confirms the independent contribution of learned memory skills. In this setting, MemSkill achieves the best average PR under both base models, with consistent gains on both Test-N and Test-C. This shows that MemSkill is not merely benefiting from demonstrations, but can provide useful experience abstraction when the agent must rely more directly on constructed memory.

Overall, these results show that MemSkill generalizes to interactive tool-use environments beyond conversational and embodied household tasks. Under the default demonstration-based setting, MemSkill remains competitive even when demonstrations provide strong task guidance, and brings clearer gains when the base model leaves more room for memory-based improvement. Under the controlled no-demonstration setting, MemSkill consistently improves pass rate over prior memory methods, further supporting the effectiveness of skill-conditioned memory construction for complex multi-step tool use.





\subsection{Experimental Results on Small Models}
\label{appendix:small_model}



We further evaluate MemSkill with Llama-3.1-8B-Instruct to examine whether the learned memory skills remain effective when the base model has more limited capacity. As shown in Table~\ref{tab:exp_small_model}, MemSkill consistently achieves the best performance across both LoCoMo and LongMemEval. Compared with prior memory methods, MemSkill improves both F1 and L-J on the two conversational benchmarks, indicating that skill-conditioned memory construction can provide useful support even for smaller backbone models.

The improvement is also maintained on LongMemEval, where no training is performed using this base model or dataset. This suggests that the learned memory skills are not tightly coupled to a specific backbone, and can transfer to smaller models under long-context conversational settings. Overall, these results show that MemSkill does not rely solely on the capability of a large base model. Instead, its learned memory skills provide complementary gains that remain effective under a more resource-efficient model setting.




\begin{table*}[h]
    \footnotesize
    \centering
    \setlength{\tabcolsep}{5.5pt}
    \caption{\textbf{Experimental results on LoCoMo and LongMemEval using Llama-3.1-8B-Instruct.}}
    \label{tab:exp_small_model}

    \begin{tabular}{>{\centering\arraybackslash}m{0.5cm} l cc cc c}
        \toprule
        \multirow{3}{*}{\textbf{Model}} &
        \multirow{3}{*}{\textbf{Methods}} &
        \multicolumn{5}{c}{\textbf{Conversational Benchmarks}} \\
        \cmidrule(lr){3-7}
        & &
        \multicolumn{2}{c}{\textbf{LoCoMo}} &
        \multicolumn{2}{c}{\textbf{$^\blacktriangle$LongMemEval}} &
        \multicolumn{1}{c}{\textbf{Avg.}} \\
        \cmidrule(lr){3-4} \cmidrule(lr){5-6} \cmidrule(lr){7-7}
        & &
        \textbf{F1} & \textbf{L-J} &
        \textbf{F1} & \textbf{L-J} &
        \textbf{L-J} \\
        \midrule

        % ===================== Base model block 1 =====================
        \multirow{8}{*}{\rotatebox[origin=c]{90}{\shortstack{\textbf{LLaMA3.1}\\\textbf{8B-Instruct}}}} &
        ReadAgent  & 17.53 & 21.49 & 9.45 & 18.80 & 20.15 \\
        & MemoryBank & 18.84 & 22.36 & 13.27 & 24.59 & 23.48 \\
        & A-MEM      & 19.95 & 25.62 & 16.12 & 25.71 & 25.67 \\
        & Mem0       & 14.26 & 17.12 & 14.59 & 26.38 & 21.75 \\
        & LangMem    & 13.36 & 15.71 & 6.52 & 17.16 & 16.44 \\
        & MemoryOS   & 17.68 & 23.08 & 7.88 & 18.24 & 20.66 \\
        & LightMem   & 19.74 & 26.85 & 17.68 & 29.40 & 28.13 \\
        & \cellcolor{hl-llama}\textbf{MemSkill} &
        \cellcolor{hl-llama}\textbf{20.59} & \cellcolor{hl-llama}\textbf{27.23} & \cellcolor{hl-llama}\textbf{18.13} & \cellcolor{hl-llama}\textbf{30.20} &
        \cellcolor{hl-llama}\textbf{28.72} \\
        \bottomrule

        \multicolumn{7}{l}{\small \textbf{Bold} indicates the best score within each base model block.} \\
        \multicolumn{7}{l}{\small $^\blacktriangle$ indicates no training using this base model or dataset.}
    \end{tabular}
\end{table*}









\subsection{Training Stability}
\label{appendix:train_stability}



MemSkill incorporates several mechanisms to improve the stability of skill learning and evolution. First, we maintain a snapshot of the skill bank after each evolution round. If the updated skill bank underperforms the best-performing snapshot observed so far, we roll back to the best snapshot and continue subsequent evolution from it. This prevents occasional harmful skill updates from permanently degrading the memory system.

Second, hard cases are selected according to difficulty scores, so the designer is guided by recurring failure patterns rather than noisy random examples. This encourages each evolution round to focus on systematic weaknesses of the current skill bank. Third, skill evolution is performed at controlled intervals, with a capped number of skill modifications in each round. This avoids abrupt changes to the skill bank and makes the learning process more gradual.

Together, these mechanisms make the evolution process less sensitive to noisy feedback and reduce the risk of unstable skill drift. In practice, we observe that the learned skill bank improves progressively over evolution rounds, suggesting that the snapshot rollback, hard-case selection, and controlled update strategy provide a stable basis for self-evolving memory skills.
More details can be found in Appendix~\ref{appendix:detail_component}







\section{More Implementation Details}
\label{appendix:imple_detail}




\subsection{Evaluation Details}
\label{appendix:exp_eval}

\textbf{LLM judge and infrastructure.}
We use \texttt{openai/gpt-oss-120b} as the LLM judge (judge prompts can be found in Appendix~\ref{appendix:prompts}). All API-based models are accessed through the NV NIM API\footnote{\url{https://docs.nvidia.com/nim/}} and Together API\footnote{\url{https://docs.together.ai/}}. Training is conducted on NVIDIA A6000 GPUs.

\textbf{LoCoMo~\citep{maharana2024evaluating-locomo}.}
LoCoMo contains 10 long interaction samples, each paired with roughly 200 training queries on average. We split the dataset by sample into train, validation, and test sets with a 6/2/2 ratio. We further remove \emph{adversarial} queries, since their supporting evidence is absent from the provided context and may introduce noisy supervision during training.

\textbf{LongMemEval~\citep{wu2024longmemeval}.}
We use the LongMemEval-S split, where each example contains an ultra-long conversation of roughly 100K tokens. We remove abstention questions, since they do not require retrieving or constructing useful memory from the conversation history. We split the remaining data into train, validation, and test sets. Although the training split is not used for learning in this transfer setting, the validation split is used to tune dataset-specific configurations. We then conduct transfer evaluation on a stratified test sample of about one-fifth of the dataset, approximately 100 samples, ensuring coverage of different question types for a comprehensive assessment.

\textbf{ALFWorld~\citep{shridhar2020alfworld}.}
We first collect expert trajectories from the training split and use them as the corpus for memory or experience construction. We then evaluate on the official ALF-Seen and ALF-Unseen splits. More training configuration details can be found in Appendix~\ref{appendix:detail_alfworld}.

\textbf{HotpotQA~\citep{yang2018hotpotqa}.}
We use HotpotQA to study transfer under distribution shift, following the evaluation protocol of~\citep{yu2025memagent}. Specifically, we evaluate on three context-length settings with increasing difficulty, corresponding to 50, 100, and 200 concatenated documents, denoted as \texttt{eval\_50}, \texttt{eval\_100}, and \texttt{eval\_200}. Unless otherwise specified, all results in this part use LLaMA as the base model and report the LLM-judge score (L-J).

\textbf{Span-level evaluation.}
During evaluation, we construct memory at the span level with a default span size of 512 tokens, rather than updating memory turn by turn. This substantially reduces the number of LLM calls and improves evaluation efficiency.









\subsection{More Details of the Designer}
\label{appendix:detail_component}

\paragraph{Hard-case buffer and representative case mining.}
The designer maintains a sliding \emph{hard-case buffer} that tracks recently challenging evaluation cases without growing unbounded. 
Each case stores the query, the retrieved memories used to answer it, the model prediction, the reference answer, the resulting task reward (e.g., F1), and a failure counter that records how many times the case has been answered incorrectly.
To prioritize cases that are both low-reward and repeatedly failed, we assign each case a difficulty score
\begin{equation}
d(q) \;=\; \big(1 - r(q)\big)\cdot c(q),
\end{equation}
where $r(q)\in[0,1]$ is the task reward for query $q$ and $c(q)$ is its cumulative failure count within the buffer window. Higher $d(q)$ indicates more critical cases that should be examined first.

To encourage coverage over \emph{diverse} failure types, we further cluster hard cases by semantic similarity of their queries and mine representative cases from each cluster.
For example, in LoCoMo, some queries focus on temporal cues (e.g., \emph{when} an event happened) while others emphasize locations (e.g., \emph{where} something occurred). 
Clustering helps separate these semantic types so the designer feedback is not dominated by a single frequent error mode, improving diversity and completeness of the mined supervision.

\paragraph{Exploration incentive for newly introduced skills.}
After each evolution round, the designer may introduce new skills that the controller has not yet learned to utilize.
To facilitate adoption, we apply a short post-update exploration phase by biasing the controller toward new skills directly at the logit level.
Let $\mathcal{S}_{\text{new}} \subseteq \mathcal{S}_t$ denote the newly added skills, and let $p_\theta(i\mid h_t)=\mathrm{softmax}(z_t)_i$ be the controller distribution at step $t$.
We encourage the total probability mass assigned to new skills to reach a target threshold $\tau_q$:
\begin{equation}
\sum_{i\in \mathcal{S}_{\text{new}}} p_\theta(i\mid h_t)\;\ge\;\tau_q,
\qquad \tau_q\in[0,1].
\label{eq:new_skill_mass_constraint}
\end{equation}
When the constraint in Eq.~\eqref{eq:new_skill_mass_constraint} is violated, we add a uniform logit gain $\delta_q$ to all new skills,
\begin{equation}
z'_{t,i} \;=\;
\begin{cases}
z_{t,i} + \delta_q, & i\in \mathcal{S}_{\text{new}},\\
z_{t,i}, & \text{otherwise},
\end{cases}
\qquad
p'_\theta(\cdot\mid h_t)=\mathrm{softmax}(z'_t),
\end{equation}
where $\delta_q$ is chosen as the minimal value that makes $\sum_{i\in \mathcal{S}_{\text{new}}} p'_\theta(i\mid h_t)\ge\tau_q$.
By operating on logits, this mechanism preserves the controller architecture and provides a smooth probability-level encouragement toward new skills.

We apply this incentive for the first $T_{\text{explore}}{=}50$ training steps after each evolution round.
To avoid persistent bias, the target threshold decays linearly within this window:
\begin{equation}
\tau_q \;=\; \tau_0 \cdot \Big(1 - \frac{q}{T_{\text{explore}}}\Big),
\qquad q=0,1,\dots,T_{\text{explore}}-1,
\end{equation}
with default $\tau_0{=}0.3$.
This schedule provides strong initial exploration and then gradually fades, yielding a smooth transition back to the controller's learned selection behavior.



\paragraph{Early stopping and rollback based on stabilized rewards.}
MemSkill performs skill evolution periodically, where each evolution cycle consists of a fixed number of controller-training steps (e.g., 100 steps) on the current skill bank. Because the reward signal can be volatile immediately after a skill-bank update, we assess whether a cycle improves performance using a \emph{stabilized} reward estimate: we compute the average task reward over the \emph{last quarter} of training steps within the cycle, and treat this value as the cycle's score.

Let $L$ denote the number of controller-training steps per cycle and $\{r_t\}_{t=1}^{L}$ the step-level rewards within the cycle. We define the cycle score as
\begin{equation}
\bar{r}_{\text{tail}} \;=\; \frac{1}{L/4}\sum_{t=3L/4+1}^{L} r_t .
\end{equation}
We compare $\bar{r}_{\text{tail}}$ against the best score observed so far. If the current cycle does not improve this criterion, then before performing the next skill evolution step, we roll back the skill bank to the previously best-performing snapshot and restart evolution from that snapshot. This rollback prevents compounding degradations from suboptimal designer updates.

Finally, if the stabilized reward fails to improve for several consecutive evolution cycles (we use a fixed patience), we early stop training and return the best skill bank snapshot encountered during training.







\subsection{Details on ALFWorld Training}
\label{appendix:detail_alfworld}

ALFWorld differs from the other benchmarks in that it is an interactive environment rather than a static text corpus. To instantiate MemSkill in this setting, we first convert ALFWorld into an offline training protocol by collecting expert trajectories on the training split. Each trajectory records the agent's interaction sequence (observations, actions, and outcomes) and serves as an interaction trace for memory construction.

\paragraph{Task-type grouping.}
ALFWorld tasks naturally fall into a small number of recurring goal templates. Following common practice, we group trajectories by task type (i.e., goal template), such as \textsc{Pick \& Place} (put an object into/on a target receptacle), \textsc{Clean \& Place} (clean an object and then place it), \textsc{Heat \& Place} (heat an object and then place it), and \textsc{Cool \& Place} (cool an object and then place it).\footnote{We use the task template provided by the environment to define task types.}

\paragraph{Experience corpus vs.\ evaluation cases.}
To fit ALFWorld into our training framework, we construct per-type train-time data splits from the offline expert trajectories. For each task type, we randomly sample a subset of trajectories as the \emph{experience corpus} used for memory construction, and sample another \emph{non-overlapping} subset of trajectories from the same type as \emph{evaluation cases}. During training, MemSkill builds a trajectory-specific memory bank from the experience corpus (span by span, via controller and executor), and then evaluates the constructed memory on the evaluation cases to obtain task reward and to log failure cases.

\paragraph{Motivation.}
Using non-overlapping trajectories from the \emph{same} task type for experience construction and evaluation provides a controlled generalization signal: trajectories within a type share goal structure and recurrent interaction patterns, making memories and skills more transferable across different instances of the same template. This setup encourages MemSkill to learn reusable memory skills that capture type-level regularities (e.g., relevant object states and action prerequisites) rather than overfitting to a single trajectory, while still ensuring that evaluation traces are held out from the traces used to build memory.









\subsection{Details on Training Objectives}
\label{appendix:rl_objective}



This part details the reinforcement learning objective used to optimize the controller in MemSkill when each decision selects an ordered Top-$K$ \emph{set} of skills without replacement.

\paragraph{Episode, states, and Top-$K$ actions.}
Training iterates over interaction traces (episodes). For a trace, MemSkill processes spans sequentially. 
At step $t$, the controller observes the raw state $s_t \triangleq (x_t,M_t)$, represented by the learned state embedding $h_t$ defined in Section~\ref{sec:controller}.
Let $\mathcal{S}_t=\{1,\dots,N_t\}$ denote the current skill bank, whose size $N_t$ may change as the designer evolves skills.
The controller computes logits $z_t\in\mathbb{R}^{N_t}$ by applying the shared scorer to all state-skill pairs, and induces
\begin{equation}
p_\theta(i\mid h_t)=\mathrm{softmax}(z_t)_i.
\end{equation}
Instead of sampling a single skill, the controller selects an \emph{ordered} Top-$K$ set
$A_t=(a_{t,1},\dots,a_{t,K})$ \emph{without replacement}, implemented via Gumbel-Top-$K$ sampling~\citep{kool2019stochastic} (i.e., adding i.i.d.\ Gumbel noise to logits and taking the top-$K$ indices).

\paragraph{Joint probability of Top-$K$ without-replacement selection.}
For PPO-style policy optimization, we need the joint probability of sampling the ordered set $A_t$ under the without-replacement process. This probability can be written as
\begin{equation}
\pi_\theta(A_t \mid h_t)
= \prod_{j=1}^{K}
\frac{p_\theta(a_{t,j}\mid h_t)}
{1-\sum_{\ell<j} p_\theta(a_{t,\ell}\mid h_t)} ,
\label{eq:topk_joint_prob_appendix}
\end{equation}
with the corresponding joint log-probability
\begin{equation}
\log \pi_\theta(A_t \mid h_t)
= \sum_{j=1}^{K}\Big(
\log p_\theta(a_{t,j}\mid h_t)
-\log\big(1-\sum_{\ell<j} p_\theta(a_{t,\ell}\mid h_t)\big)
\Big).
\label{eq:topk_joint_logprob_appendix}
\end{equation}
When $K=1$, Eq.~\ref{eq:topk_joint_prob_appendix} reduces to the standard single-action case.

\paragraph{Rewards from memory-dependent evaluation.}
For each trace, after processing all spans and constructing the trace-specific memory bank, we evaluate the memory bank on the trace's memory-dependent training queries and obtain a scalar task score (e.g., F1 or success rate). We treat this score as the episode-level reward:
\begin{equation}
R \triangleq \mathrm{Eval}(\text{memory bank}; \text{training queries}) \in \mathbb{R}.
\end{equation}
This reward is then assigned to the sequence of controller decisions within the trace. Concretely, we use standard return computation with discount factor $\gamma$:
\begin{equation}
G_t=\sum_{\tau=t}^{T}\gamma^{\tau-t}r_\tau,
\end{equation}
where $r_\tau$ is the per-step reward. In our default setting, reward is provided only after memory construction completes, i.e., $r_T=R$ and $r_\tau=0$ for $\tau<T$, so $G_t=\gamma^{T-t}R$.
We learn a value function $V_\phi(h_t)$ and compute advantages $\hat{A}_t$ using generalized advantage estimation (GAE).

\paragraph{PPO objective with Top-$K$ actions.}
We optimize the controller using proximal policy optimization (PPO)~\citep{schulman2017proximal}, replacing the standard single-action log-probability with the Top-$K$ joint log-probability in Eq.~\ref{eq:topk_joint_logprob_appendix}.
Let $\theta_{\text{old}}$ denote the parameters of the behavior policy used to collect rollouts.
For simplicity, we use $h_t$ to denote the state representation computed under the policy being evaluated.
Define the importance ratio:
\begin{equation}
r_t(\theta)=
\frac{\pi_\theta(A_t\mid h_t)}{\pi_{\theta_{\text{old}}}(A_t\mid h_t)}
=
\exp\Big(\log\pi_\theta(A_t\mid h_t)-\log\pi_{\theta_{\text{old}}}(A_t\mid h_t)\Big).
\label{eq:topk_ratio}
\end{equation}
The clipped surrogate policy objective is
\begin{equation}
\mathcal{L}_{\text{policy}}(\theta)
=
\mathbb{E}_t\Big[
\min\big(r_t(\theta)\,\hat{A}_t,\;
\mathrm{clip}(r_t(\theta),1-\epsilon,1+\epsilon)\,\hat{A}_t\big)
\Big].
\label{eq:ppo_policy_topk}
\end{equation}
We additionally optimize a value function and include an entropy bonus for exploration:
\begin{align}
\mathcal{L}_{\text{value}}(\phi)
&=
\mathbb{E}_t\Big[
\big(V_\phi(h_t)-G_t\big)^2
\Big], \label{eq:ppo_value}\\
\mathcal{H}(\theta)
&=
\mathbb{E}_t\big[H(p_\theta(\cdot\mid h_t))\big], \label{eq:ppo_entropy}
\end{align}
where $H(\cdot)$ is the entropy of the categorical distribution over all skills.
The overall objective (to maximize) is
\begin{equation}
\max_{\theta,\phi}\;
\mathcal{L}_{\text{policy}}(\theta)
- c_v\,\mathcal{L}_{\text{value}}(\phi)
+ c_H\,\mathcal{H}(\theta).
\label{eq:ppo_full_topk}
\end{equation}
In implementation, we minimize the negative of Eq.~\ref{eq:ppo_full_topk}.

\paragraph{Gumbel-Top-$K$ exploration.}
To sample Top-$K$ skills without replacement during rollout collection, we use Gumbel-Top-$K$ sampling: at each step we draw i.i.d.\ Gumbel noise $\{g_i\}_{i=1}^{N_t}$, form perturbed logits $\tilde{z}_{t,i}=z_{t,i}+g_i$, and take the indices of the $K$ largest $\tilde{z}_{t,i}$ to obtain $A_t$.
This provides stochastic exploration over skill subsets while remaining compatible with PPO through the joint probability in Eq.~\ref{eq:topk_joint_prob_appendix}.
For training stability, entropy regularization is computed from the base categorical distribution $p_\theta(\cdot\mid h_t)$ over all skills (Eq.~\ref{eq:ppo_entropy}), which encourages exploration of the evolving skill bank even though the executed action is a Top-$K$ set.










\section{Case Study}
\label{appendix:case_study}





\subsection{Initial Primitive Skills}


\begin{center}
\begin{tcolorbox}[title={Initial Primitive Skill - INSERT}]
{
Skill: Insert New Memory \\\\
Description: Memory management skill for capturing new, durable facts from the current text chunk that are not already in memory. \\\\
Purpose: Capture new, durable facts from the current text chunk that are missing in memory. \\\\
When to use: \\
- The text chunk introduces new facts, events, plans, or context worth storing. \\
- The information is stable and likely useful later. \\\\
How to apply: \\
- Compare against retrieved memories to avoid duplicates. \\
- Split distinct facts into separate items. \\
- Keep each item concise and specific. \\\\
Constraints: \\
- Skip trivial, fleeting, or speculative content. \\
- Do not update or delete existing memories. \\\\
Action type: INSERT only.
}
\label{appendix:case_insert}
\end{tcolorbox}
\end{center}



\begin{center}
\begin{tcolorbox}[title={Initial Primitive Skill - UPDATE}]
{
Skill: Update Existing Memory \\\\
Description: Memory management skill for revising an existing memory item when the text chunk provides corrections or new details. \\\\
Purpose: Revise a retrieved memory with new or corrected information from the text chunk. \\\\
When to use: \\
- The text chunk clarifies, corrects, or extends a retrieved memory. \\\\
How to apply: \\
- Select the best matching memory item. \\
- Merge new details into a single updated item. \\
- Preserve accurate details that still hold. \\\\
Constraints: \\
- Do not create new memories. \\
- Do not delete items. \\\\
Action type: UPDATE only.
}
\label{appendix:case_update}
\end{tcolorbox}
\end{center}





\begin{center}
\begin{tcolorbox}[title={Initial Primitive Skill - DELETE}]
{
Skill: Delete Invalid Memory \\\\
Description: Memory management skill for removing memory items that are incorrect, outdated, or superseded. \\\\
Purpose: Remove a retrieved memory that is wrong, outdated, or superseded by the text chunk. \\\\
When to use: \\
- The text chunk clearly contradicts a memory. \\
- A plan or fact is explicitly canceled or replaced. \\\\
How to apply: \\
- Only delete when evidence is explicit. \\\\
Constraints: \\
- If uncertain, prefer no action over deletion. \\\\
Action type: DELETE only.
}
\label{appendix:case_delete}
\end{tcolorbox}
\end{center}




\begin{center}
\begin{tcolorbox}[title={Initial Primitive Skill - SKIP}]
{
Skill: No Operation \\\\
Description: Memory management skill for confirming that no memory changes are required. \\\\
Purpose: Confirm no memory changes are needed for the text chunk. \\\\
When to use: \\
- The text chunk contains no new, corrective, or actionable information. \\\\
Constraints: \\
- Emit NOOP only if none of the selected skills produce actions. \\\\
Action type: NOOP only.
}
\label{appendix:case_noop}
\end{tcolorbox}
\end{center}





\subsection{Evolved Skills on LoCoMo}



% ---------- capture_activity_details ----------
\begin{center}
\begin{tcolorbox}[
  title={Evolved Skill on LoCoMo - CAPTURE\_ACTIVITY\_DETAILS},
  colback=blue!2!white,
  colframe=blue!35!black,
  colbacktitle=blue!15!white,
  coltitle=black,
  coltext=black
]
{
Skill: Capture Activity Details \\\\
Purpose: Capture detailed information about activities mentioned in the text chunk, including the type of activity, location, participants, temporal details, and any relevant contextual information. \\\\
When to use: \\
- The text chunk mentions a specific activity or event with contextual details. \\\\
How to apply: \\
- Identify the key elements of the activity (e.g., type, location, participants, temporal details). \\
- Capture any relevant contextual information that provides additional insight into the activity. \\
- Keep the activity details specific, actionable, and concise. \\\\
Constraints: \\
- Focus on explicit activity details and contextual information mentioned in the text chunk. \\
- Avoid inferring activity details or context not directly stated. \\\\
Action type: INSERT only.
}
\label{appendix:case_capture_activity_details}
\end{tcolorbox}
\end{center}

% ---------- capture_entity_nuances ----------
\begin{center}
\begin{tcolorbox}[
  title={Evolved Skill on LoCoMo - CAPTURE\_ENTITY\_NUANCES},
  colback=blue!2!white,
  colframe=blue!35!black,
  colbacktitle=blue!15!white,
  coltitle=black,
  coltext=black
]
{
Skill: Capture Entity Nuances \\\\
Purpose: Capture nuanced details about entities mentioned in the text chunk, such as nicknames, aliases, or comparative statements. \\\\
When to use: \\
- The text chunk mentions an entity with nuanced details (e.g., nickname, alias, comparison). \\\\
How to apply: \\
- Identify the entity and its associated nuances. \\
- Capture these nuances in a way that distinguishes them from the entity's primary information. \\\\
Constraints: \\
- Focus on explicit nuances mentioned in the text chunk. \\
- Avoid inferring nuances not directly stated. \\\\
Action type: INSERT only.
}
\label{appendix:case_capture_entity_nuances}
\end{tcolorbox}
\end{center}

% ---------- capture_temporal_context ----------
\begin{center}
\begin{tcolorbox}[
  title={Evolved Skill on LoCoMo - CAPTURE\_TEMPORAL\_CONTEXT},
  colback=blue!2!white,
  colframe=blue!35!black,
  colbacktitle=blue!15!white,
  coltitle=black,
  coltext=black
]
{
Skill: Capture Temporal Context \\\\
Purpose: Capture the temporal context of events, activities, or facts mentioned in the text chunk, including any relevant dates, times, durations, or sequential information. \\\\
When to use: \\
- The text chunk mentions a specific event, activity, or fact with associated temporal information. \\\\
How to apply: \\
- Identify the key temporal elements (e.g., start time, end time, duration, sequence). \\
- Capture the temporal context in a concise and specific format, considering any sequential relationships. \\\\
Constraints: \\
- Focus on explicit temporal information mentioned in the text chunk. \\
- Avoid inferring temporal details not directly stated. \\\\
Action type: INSERT only.
}
\label{appendix:case_capture_temporal_context}
\end{tcolorbox}
\end{center}

% ---------- delete ----------
\begin{center}
\begin{tcolorbox}[
  title={Evolved Skill on LoCoMo - DELETE},
  colback=blue!2!white,
  colframe=blue!35!black,
  colbacktitle=blue!15!white,
  coltitle=black,
  coltext=black
]
{
Skill: Delete Invalid Memory \\\\
Purpose: Remove a retrieved memory that is wrong, outdated, or superseded by the text chunk. \\\\
When to use: \\
- The text chunk clearly contradicts a memory. \\
- A plan or fact is explicitly canceled or replaced. \\\\
How to apply: \\
- Only delete when evidence is explicit. \\\\
Constraints: \\
- If uncertain, prefer no action over deletion. \\\\
Action type: DELETE only.
}
\label{appendix:case_delete}
\end{tcolorbox}
\end{center}

% ---------- handle_entity_relationships ----------
\begin{center}
\begin{tcolorbox}[
  title={Evolved Skill on LoCoMo - HANDLE\_ENTITY\_RELATIONSHIPS},
  colback=blue!2!white,
  colframe=blue!35!black,
  colbacktitle=blue!15!white,
  coltitle=black,
  coltext=black
]
{
Skill: Handle Entity Relationships \\\\
Purpose: Capture and manage complex relationships between entities mentioned in the text chunk, including nuanced details. \\\\
When to use: \\
- The text chunk mentions interactions, associations, or relationships between entities with specific details. \\\\
How to apply: \\
- Identify the entities involved and their roles in the relationship. \\
- Capture the nature of the relationship and any nuanced details (e.g., nicknames, comparative statements). \\
- Update existing memories to reflect the new relationship information. \\\\
Constraints: \\
- Focus on explicit relationships mentioned in the text chunk. \\
- Avoid inferring relationships not directly stated. \\\\
Action type: INSERT only.
}
\label{appendix:case_handle_entity_relationships}
\end{tcolorbox}
\end{center}

% ---------- insert ----------
\begin{center}
\begin{tcolorbox}[
  title={Evolved Skill on LoCoMo - INSERT},
  colback=blue!2!white,
  colframe=blue!35!black,
  colbacktitle=blue!15!white,
  coltitle=black,
  coltext=black
]
{
Skill: Insert New Memory \\\\
Purpose: Capture new, durable facts from the current text chunk that are missing in memory, including specific temporal details such as dates or time frames and detailed activity information. \\\\
When to use: \\
- The text chunk introduces new facts, events, plans, or context worth storing. \\
- The information is stable and likely useful later. \\\\
How to apply: \\
- Compare against retrieved memories to avoid duplicates. \\
- Split distinct facts into separate items. \\
- Keep each item concise and specific, including relevant temporal information and activity details. \\\\
Constraints: \\
- Skip trivial, fleeting, or speculative content. \\
- Do not update or delete existing memories. \\\\
Action type: INSERT only.
}
\label{appendix:case_insert}
\end{tcolorbox}
\end{center}

% ---------- noop ----------
\begin{center}
\begin{tcolorbox}[
  title={Evolved Skill on LoCoMo - NOOP},
  colback=blue!2!white,
  colframe=blue!35!black,
  colbacktitle=blue!15!white,
  coltitle=black,
  coltext=black
]
{
Skill: No Operation \\\\
Purpose: Confirm no memory changes are needed for the text chunk. \\\\
When to use: \\
- The text chunk contains no new, corrective, or actionable information. \\\\
Constraints: \\
- Emit NOOP only if none of the selected skills produce actions. \\\\
Action type: NOOP only.
}
\label{appendix:case_noop}
\end{tcolorbox}
\end{center}

% ---------- refine_temporal_details_with_context ----------
\begin{center}
\begin{tcolorbox}[
  title={Evolved Skill on LoCoMo - REFINE\_TEMPORAL\_DETAILS\_WITH\_CONTEXT},
  colback=blue!2!white,
  colframe=blue!35!black,
  colbacktitle=blue!15!white,
  coltitle=black,
  coltext=black
]
{
Skill: Refine Temporal Details with Context \\\\
Purpose: Update the temporal context of existing memories with new information from the text chunk, considering the context in which the information is provided. \\\\
When to use: \\
- The text chunk provides new or corrected temporal information relevant to an existing memory, and the context suggests a need for refinement. \\\\
How to apply: \\
- Identify the relevant existing memory and its current temporal context. \\
- Update the temporal details to reflect the new information, ensuring consistency with the provided context. \\\\
Constraints: \\
- Focus on explicit temporal information mentioned in the text chunk and supported by the context. \\
- Avoid inferring temporal details not directly stated or implied by the context. \\\\
Action type: UPDATE only.
}
\label{appendix:case_refine_temporal_details_with_context}
\end{tcolorbox}
\end{center}

% ---------- update ----------
\begin{center}
\begin{tcolorbox}[
  title={Evolved Skill on LoCoMo - UPDATE},
  colback=blue!2!white,
  colframe=blue!35!black,
  colbacktitle=blue!15!white,
  coltitle=black,
  coltext=black
]
{
Skill: Update Existing Memory \\\\
Purpose: Revise a retrieved memory with new or corrected information from the text chunk, including entity-specific details. \\\\
When to use: \\
- The text chunk clarifies, corrects, or extends a retrieved memory. \\
- The text chunk provides new information about a specific entity or its activities. \\\\
How to apply: \\
- Select the best matching memory item. \\
- Merge new details into a single updated item. \\
- Preserve accurate details that still hold, and ensure entity-specific information is accurately captured and updated. \\\\
Constraints: \\
- Do not create new memories. \\
- Do not delete items. \\\\
Action type: UPDATE only.
}
\label{appendix:case_update}
\end{tcolorbox}
\end{center}












\subsection{Evolved Skills on ALFWorld}


% ---------- capture_action_constraints ----------
\begin{center}
\begin{tcolorbox}[
  title={Evolved Skill on ALFWorld - CAPTURE\_ACTION\_CONSTRAINTS},
  colback=orange!6,
  colframe=orange!75!black,
  coltitle=black
]
{
Skill: Capture Action Constraints \\\\
Purpose: Capture detailed constraints on actions, including object states and movements, necessary for task completion. \\\\
When to use: \\
- The text chunk mentions constraints on actions, including object states and movements. \\
- The constraints are crucial for future task steps. \\\\
How to apply: \\
- Identify the action, its constraints, and relevant object states and movements from the text chunk. \\
- Create a new memory item with the action-constraint pair, including object states and movements. \\\\
Constraints: \\
- Only capture constraints on actions relevant to the task. \\
- Update existing constraint memories if new information is provided. \\\\
Action type: INSERT only.
}
\label{appendix:case_capture_action_constraints}
\end{tcolorbox}
\end{center}

% ---------- delete ----------
\begin{center}
\begin{tcolorbox}[
  title={Evolved Skill on ALFWorld - DELETE},
  colback=orange!6,
  colframe=orange!75!black,
  coltitle=black
]
{
Skill: Delete Invalid Memory \\\\
Purpose: Remove a retrieved memory that is wrong, outdated, or superseded by the text chunk. \\\\
When to use: \\
- The text chunk clearly contradicts a memory. \\
- A plan or fact is explicitly canceled or replaced. \\\\
How to apply: \\
- Only delete when evidence is explicit. \\\\
Constraints: \\
- If uncertain, prefer no action over deletion. \\\\
Action type: DELETE only.
}
\label{appendix:case_delete}
\end{tcolorbox}
\end{center}

% ---------- insert ----------
\begin{center}
\begin{tcolorbox}[
  title={Evolved Skill on ALFWorld - INSERT},
  colback=orange!6,
  colframe=orange!75!black,
  coltitle=black
]
{
Skill: Insert New Memory \\\\
Purpose: Capture new, durable facts, including procedural knowledge and action sequences, from the current text chunk that are missing in memory. \\\\
When to use: \\
- The text chunk introduces new facts, events, plans, or context worth storing. \\
- The information is stable and likely useful later. \\\\
How to apply: \\
- Compare against retrieved memories to avoid duplicates. \\
- Split distinct facts into separate items, including action sequences. \\
- Keep each item concise and specific. \\\\
Constraints: \\
- Skip trivial, fleeting, or speculative content. \\
- Do not update or delete existing memories. \\\\
Action type: INSERT only.
}
\label{appendix:case_insert}
\end{tcolorbox}
\end{center}

% ---------- noop ----------
\begin{center}
\begin{tcolorbox}[
  title={Evolved Skill on ALFWorld - NOOP},
  colback=orange!6,
  colframe=orange!75!black,
  coltitle=black
]
{
Skill: No Operation \\\\
Purpose: Confirm no memory changes are needed for the text chunk. \\\\
When to use: \\
- The text chunk contains no new, corrective, or actionable information. \\\\
Constraints: \\
- Emit NOOP only if none of the selected skills produce actions. \\\\
Action type: NOOP only.
}
\label{appendix:case_noop}
\end{tcolorbox}
\end{center}

% ---------- track_object_location ----------
\begin{center}
\begin{tcolorbox}[
  title={Evolved Skill on ALFWorld - TRACK\_OBJECT\_LOCATION},
  colback=orange!6,
  colframe=orange!75!black,
  coltitle=black
]
{
Skill: Track Object Location \\\\
Purpose: Explicitly track the location and state of an object necessary for task completion. \\\\
When to use: \\
- The text chunk mentions an object's location or state. \\
- The object's location or state is crucial for future task steps. \\\\
How to apply: \\
- Identify the object, its location, and relevant state from the text chunk. \\
- Create a new memory item with the object-location-state triplet. \\\\
Constraints: \\
- Only track locations and states of objects relevant to the task. \\
- Update existing location memories if new information is provided. \\\\
Action type: INSERT only.
}
\label{appendix:case_track_object_location}
\end{tcolorbox}
\end{center}

% ---------- track_object_movements ----------
\begin{center}
\begin{tcolorbox}[
  title={Evolved Skill on ALFWorld - TRACK\_OBJECT\_MOVEMENTS},
  colback=orange!6,
  colframe=orange!75!black,
  coltitle=black
]
{
Skill: Track Object Movements \\\\
Purpose: Track movements of objects necessary for task completion. \\\\
When to use: \\
- The text chunk mentions an object's movement. \\
- The object's movement is crucial for future task steps. \\\\
How to apply: \\
- Identify the object and its movement from the text chunk. \\
- Create a new memory item with the object-movement pair. \\\\
Constraints: \\
- Only track movements of objects relevant to the task. \\
- Update existing movement memories if new information is provided. \\\\
Action type: INSERT only.
}
\label{appendix:case_track_object_movements}
\end{tcolorbox}
\end{center}

% ---------- update ----------
\begin{center}
\begin{tcolorbox}[
  title={Evolved Skill on ALFWorld - UPDATE},
  colback=orange!6,
  colframe=orange!75!black,
  coltitle=black
]
{
Skill: Update Existing Memory \\\\
Purpose: Revise a retrieved memory with new or corrected information from the text chunk. \\\\
When to use: \\
- The text chunk clarifies, corrects, or extends a retrieved memory. \\\\
How to apply: \\
- Select the best matching memory item. \\
- Merge new details into a single updated item. \\
- Preserve accurate details that still hold. \\\\
Constraints: \\
- Do not create new memories. \\
- Do not delete items. \\\\
Action type: UPDATE only.
}
\label{appendix:case_update}
\end{tcolorbox}
\end{center}










\section{Prompts}
\label{appendix:prompts}






\begin{center}
\begin{tcolorbox}[title={LoCoMo Answer Prompt}]
{
Based on the above context, write an answer in the form of a short phrase for the following question. Answer with exact words from the context whenever possible. \\

Question: \{\} Short answer:
}
\label{appendix:case_locomo}
\end{tcolorbox}
\end{center}





\begin{center}
\begin{tcolorbox}[title={LongMemEval Answer Prompt}]
{
I will give you several history chats between you and a user. Please answer the question based on the relevant chat history. \\
History Chats: \{\} \\
Current Date: \{\} \\
Question: \{\} \\
Short Answer:
}
\label{appendix:case_longmemeval}
\end{tcolorbox}
\end{center}




\begin{center}
\begin{tcolorbox}[title={HotpotQA Answer Prompt}]
{
Based on the following context, answer the question. The question may require reasoning across multiple pieces of information. \\
\{context\}

Question: \{question\} \\

Instructions: \\
- Read the context carefully and identify relevant information \\
- If the answer can be found in the context, provide a short, precise answer \\
- Output your answer within \textless answer\textgreater \textless /answer\textgreater tags \\

\textless answer\textgreater your answer here\textless /answer\textgreater
}
\label{appendix:case_hp}
\end{tcolorbox}
\end{center}




\begin{center}
\begin{tcolorbox}[title={ALFWorld Env Interaction Prompt}]
{
You are controlling a text-based ALFWorld environment. Your job: choose the NEXT action as ONE text command. Output ONLY the command string, with no extra text. You MUST choose an action from the admissible actions list and copy it EXACTLY. \\\\
Goal: \{goal\} \\
Retrieved procedural tips (optional, short \& actionable): \{retrieved\_memory\} \\
Interaction history so far (most recent info matters most): \{history\} \\
Admissible actions (choose exactly ONE and copy it verbatim): \{actions\} \\\\
Now output exactly one line: the chosen action (must match one item above).
}
\label{appendix:case_alfworld}
\end{tcolorbox}
\end{center}








\begin{center}
\begin{tcolorbox}[title={Executor Prompt}]
{
You are a memory management executor. Apply the selected skills to the input text chunk and retrieved memories, then output memory actions. \\\\
Input Text Chunk: \{session\_text\} \\
Retrieved Memories (0-based index): \{mem\_text\} \\
Selected Skills: \{skills\_text\} \\\\
Guidelines: \\
- Apply any skill as needed; a skill may be used multiple times. \\
- Read the input text chunk carefully line by line and apply any skill as needed. \\
- Only use action types supported by the selected skills. \\
- MEMORY\_INDEX is 0-based and must reference the retrieved memories list. \\
- Output only action blocks in the format below. \\
- Do not include explanations or REASONING lines. \\\\
Output format (repeat as needed). Use ONE block per action and separate blocks with a blank line: \\\\
INSERT block: \\
ACTION: INSERT \\
MEMORY\_ITEM: [concise but complete summary with essential details] \\\\
UPDATE block: \\
ACTION: UPDATE \\
MEMORY\_INDEX: [0-based index] \\
UPDATED\_MEMORY: [concise but complete merged summary with essential updates] \\\\
DELETE block: \\
ACTION: DELETE \\
MEMORY\_INDEX: [0-based index]
}
\label{appendix:case_executor}
\end{tcolorbox}
\end{center}





\begin{center}
\begin{tcolorbox}[title={Designer Analysis Prompt}]
{
You are an expert analyst for a memory-augmented QA system. Analyze the failure cases below to identify why the system failed and how the memory management skills should change. \\

\#\# How This System Works \\
1. **Memory Storage**: The system applies memory management skills to decide what information to store from the text chunk. \\
2. **Memory Retrieval**: At question time, it retrieves the most relevant memories by semantic similarity. \\
3. **Answer Generation**: An LLM answers using the retrieved memories. \\

Failures can occur at any stage: \\
- **Storage failure**: Important information was never stored (skill missing or misapplied) \\
- **Retrieval failure**: Relevant memory exists but was not retrieved (embedding mismatch) \\
- **Memory quality failure**: Memory exists but is too vague or incomplete to answer \\

\#\# Current Memory Management Skills \\
\{operation\_bank\_description\} \\

\#\# Operation Evolution Feedback \\
\{evolution\_feedback\} \\

\#\# Failure Cases (\{num\_failure\_cases\} cases) \\
\{failure\_cases\_details\} \\

\#\# Analysis Instructions \\
This is round 1 of a reflection loop. Produce a strong initial analysis that can be critiqued and improved. \\
1. For each case, check whether the retrieved memories contain the answer or the needed evidence. \\
2. If missing, decide whether it was never stored (storage failure) or stored but too weak (memory quality failure). \\
3. If the answer is present but not retrieved, label it retrieval failure and avoid changing skills unless the pattern repeats. \\
4. Group cases into patterns tied to information types, entities, temporal details, or constraints. \\
5. For each pattern, propose a concrete skill change: add a new skill or refine an existing one to capture missing details. \\
6. Provide up to \{max\_changes\} recommendations total (use fewer if only one change is justified). \\

\#\# Output Format \\
Provide your analysis as JSON: \\
\{
    ``failure\_patterns": [
        \{
            ``pattern\_name": ``[descriptive name for this failure pattern]",
            ``affected\_cases": [list of case numbers, e.g., 1, 3, 5],
            ``root\_cause": ``[storage\_failure|retrieval\_failure|memory\_quality\_failure]",
            ``explanation": ``[why this pattern of failures is occurring]",
            ``potential\_fix": ``[what kind of operation change could address this]"
        \}
    ],
    ``recommendations": [
        \{
            ``action": ``[add\_new\_operation|refine\_existing\_operation|no\_change]",
            ``target\_operation": ``[operation name to refine, or null if adding new]",
            ``rationale": ``[clear explanation of why this is the best improvement]",
            ``priority": ``[high|medium|low]"
        \}
    ],
    ``summary": ``[1-2 sentence summary of main findings]"
\} \\

Focus on actionable insights. What specific change to the skill bank would prevent these failures?

Output ONLY the JSON, no other text.
}
\label{appendix:case_designer_1}
\end{tcolorbox}
\end{center}




\begin{center}
\begin{tcolorbox}[title={Designer Refinement Prompt (1/2)}]
{
Based on the failure analysis, propose a specific improvement to the memory operation bank. \\

\#\# Failure Analysis (from Stage 1) \\
\{analysis\_feedback\} \\

\#\# Current Operation Bank \\
\{operation\_bank\_full\} \\

\{evolution\_feedback\} \\

\#\# Your Task \\
Propose up to \{max\_changes\} improvements based on the analysis: \\

**Option A - Add New Operation**: Create a new operation if the analysis shows a capability gap (e.g., certain information types are not being captured). \\
**Option B - Refine Existing Operation**: Improve an existing operation's instruction template if the analysis shows it's not working well (e.g., memories are too vague, missing key details). \\
**Option C - No Change**: If the failures are due to retrieval issues (not operation issues), or if the current operations are already well-designed. \\

\#\# CRITICAL Requirements \\
1. instruction\_template MUST be a skill-style guide and MUST NOT include context placeholders (the executor injects the text chunk and retrieved memories) \\
2. instruction\_template MUST clearly state purpose, when to use, and constraints \\
3. instruction\_template MUST specify the allowed action type (INSERT or UPDATE only) \\
4. For new operations, `update\_type` must be either ``insert" or ``update" (delete and noop operations are not evolved at this time) \\
5. Only propose operations with update\_type ``insert" or ``update" \\
6. Avoid labels like ``ENHANCED", ``ADVANCED", or other marketing adjectives in descriptions or templates; keep phrasing neutral and task-specific \\
7. Do NOT embed output blocks; the executor handles output formatting and can apply the skill multiple times \\
8. The number of changes in the list MUST be less than {max\_changes} \\
9. Do NOT modify the same operation more than once in a single response, and do NOT refine an operation you add in the same response \\

\#\# Example of a Well-Designed Insert Operation \\
\{
    ``name": ``extract\_personal\_preferences",
    ``description": ``Memory management skill for capturing personal preferences and habits mentioned in the text chunk.",
    ``update\_type": ``insert",
    ``instruction\_template": ``Skill: Insert Preferences. Purpose: Capture personal preferences, habits, or opinions stated in the text chunk. When to use: The text chunk mentions likes, dislikes, routines, or goals tied to a person. How to apply: Attribute the preference to the correct person. Keep the preference specific and actionable. Constraints: Avoid one-off or ambiguous statements. Action type: INSERT only."
\} \\


\#\# Output Format \\
Respond with ONE of these JSON structures:

\#\#\# One or more changes (up to {max\_changes}):
\{
    ``action": ``apply\_changes",
    ``summary": ``[overall rationale for the set of changes]",
    ``changes": [
        \{
            ``action": ``add\_new",
            ``new\_operation": \{
                ``name": ``[snake\_case\_name]",
                ``description": ``[what it does and when to use it]",
                ``instruction\_template": ``[skill-style instruction template]",
                ``update\_type": ``[insert|update]",
                ``reasoning": ``[how this addresses the identified failures]"
            \}
        \},
        \{
            ``action": ``refine\_existing",
            ``refined\_operation": \{
                ``name": ``[existing\_operation\_name]",
                ``changes": \{
                    ``description": ``[improved description]",
                    ``instruction\_template": ``[improved template]"
                \},
                ``reasoning": ``[how these changes address the identified failures]"
            \}
        \}
    ]
\}


}
\label{appendix:case_designer_2_1}
\end{tcolorbox}
\end{center}



\begin{center}
\begin{tcolorbox}[title={Designer Refinement Prompt (2/2)}]
{
\#\#\# No changes needed: \\
\{
    ``action": ``no\_change",
    ``reasoning": ``[why the current operations are sufficient]"
\} \\

\#\# Instruction Template Structure \\
When writing instruction templates, follow this structure: \\


Skill: [Short skill name] \\
Purpose: [What this skill does] \\
When to use: \\
- [Trigger 1] \\
- [Trigger 2] \\
How to apply: \\
- [Step 1] \\
- [Step 2] \\
Constraints: \\
- [What to avoid] \\
Action type: [INSERT only | UPDATE only] \\

Output ONLY the JSON, no other text.
}
\label{appendix:case_designer_2_2}
\end{tcolorbox}
\end{center}




\begin{center}
\begin{tcolorbox}[title={LLM Judge Prompt}]
{
You are an expert judge evaluating the quality of an answer for a QA task.
Your goal is to determine whether the model's answer correctly and sufficiently
answers the given question. \\

Read the following information carefully: \\

[Question] \\
\{question\} \\

[Ground Truth Answers] \\
\{ground\_truth\} \\

[Model Answer] \\
\{model\_answer\} \\

Your evaluation criteria: \\
1. Correctness: \\
   - Is the model answer factually consistent with ANY of the correct answers? \\
   - Does it avoid contradictions or introducing false information? \\

2. Relevance: \\
   - Does the answer address the question directly without unnecessary content? \\

3. Completeness: \\
   - Does the answer include all essential information needed to fully answer the question? \\
   - Partial answers are allowed but should receive lower scores. \\

Scoring Rules: \\
- Score = 1.0 if the answer is fully correct. \\
- Score = 0.5 if the answer is partially correct but incomplete or slightly inaccurate. \\
- Score = 0.0 if the answer is incorrect, irrelevant, or contradicts the ground truth. \\

Output Format (STRICT): \\
Return your output as a JSON dictionary with two fields: \\
\{
    ``explanation": ``\textless brief explanation of your reasoning\textgreater",
    "score": \textless0.0 | 0.5 | 1.0\textgreater
\} \\

Be concise and objective. Do not include anything outside the JSON.
}
\label{appendix:case_judge}
\end{tcolorbox}
\end{center}





\section{Use of Large Language Models}
\label{app:llm_usage}

LLMs are used as part of our method and evaluation pipeline. Specifically, MemSkill uses an LLM executor for skill-conditioned memory generation, an LLM designer for skill evolution, and an LLM judge for evaluation on tasks where automatic metrics are insufficient. The corresponding prompts are provided in Appendix~\ref{appendix:prompts}.

LLM-based tools may also be used for grammar polishing and code assistance. The authors take full responsibility for all scientific claims, experimental results, figures, and references in this paper.



\section{Limitations and Societal Impact}
\label{app:limitations_impact}

\paragraph{Limitations.}
This work focuses on benchmarked research settings for skill-conditioned memory construction. Practical deployments of agent memory systems may require additional mechanisms for privacy protection, user consent, access control, and data retention. These considerations are orthogonal to the core algorithmic contribution of MemSkill and are left for future deployment-oriented studies.


\paragraph{Societal impact.}
MemSkill contributes toward more adaptive and reusable memory mechanisms for long-horizon LLM agents. Instead of repeatedly designing task-specific memory pipelines, the proposed skill-based formulation allows memory behaviors to be learned, reused, and evolved from interaction traces, which may lower the engineering barrier for building memory-augmented agents. This can benefit research and practical applications that require long-term context management, such as personalized assistants, educational agents, research support tools, and embodied task-solving systems. Since agent memory may involve user-provided or task-specific information, practical deployments should follow standard data governance practices, including user consent, privacy protection, data retention control, and mechanisms for inspecting or deleting stored memories. Our work focuses on benchmarked research settings and does not target surveillance, deception, or high-stakes decision making.